\begin{center}
{\Large \bfseries \textcolor{lyred}{第十二章 无穷级数}}
\end{center}

\vspace{6pt}
{\large \bfseries \textcolor{lyred}{第12.1 节 常数项级数的概念和性质}}

\vspace{4pt}
{\bfseries \textcolor{lyred}{一.常数项级数的概念}}

给定\{\}
n
u
,表达式
n
u
u
u
+
+
+
+
\Lambda 
\Lambda 称为\textcolor{lyred}{常数项无穷级数},简称为\textcolor{lyred}{常数项级数}, 
记为
n
n
u
\infty 
=\sum 
,它是一个形式和,
n
u 称为\textcolor{lyred}{一般项}; 
记
n
n
k
k
s
u
=
=\sum 
称为
n
n
u
\infty 
=\sum 
的\textcolor{lyred}{部分和},数列\{\}
ns
称为\textcolor{lyred}{部分和数列}; 
若极限lim
n
n
s
s
\to \infty 
=
存在,则称级数
n
n
u
\infty 
=\sum 
\textcolor{lyred}{收敛}, s 为\textcolor{lyred}{和},也记
n
s
u
u
u
=
+
+
+
+
\Lambda 
\Lambda , 
称
n
n
n
n
r
s
s
u
u
+
+
=
-
=
+
+\Lambda 为\textcolor{lyred}{余项},若lim
n
n
s
\to \infty 
不存在,则称
n
n
u
\infty 
=\sum 
\textcolor{lyred}{发散}. 
\textcolor{lyred}{注}.若lim
n
n
s
\to \infty 
=\infty ,则记
n
n
u
\infty 
=
=\infty 
\sum 
. 
\textcolor{lyred}{例}.考虑\textcolor{lyred}{几何级数},或\textcolor{lyred}{等比级数},
k
k
aq
a
aq
aq
\infty 
=
=
+
+
+
\sum 
\Lambda ,其中
a \neq 
,\textcolor{lyred}{ }
若
q \neq ,则
(
)
n
n
k
n
n
k
a
q
s
aq
a
aq
aq
aq
q
-
-
=
-
=
=
+
+
+
+
=
-
\sum 
\Lambda 
; 
(1)当
q <时, lim
n
n
q
\to \infty 
=
,故lim
n
n
a
s
q
\to \infty 
=-
,级数收敛,和为
a
s
q
=-
; 
(2)当
q >时, lim
n
n
q
\to \infty 
=\infty ,故lim
n
n
s
\to \infty 
=\infty ,级数发散; 
(3)当
q =时,
q =\pm ,显然lim
n
n
s
\to \infty 
不存在,级数发散; 
因此,当
q <时,
k
k
aq
\infty 
=\sum 
收敛;当
q \ge 时,
k
k
aq
\infty 
=\sum 
发散. 
\textcolor{lyred}{例}.讨论
(
)
1 2
2 3
3 4
n
n n
\infty 
=
=
+
+
+
+
\cdots 
\cdots 
\cdots 
\sum 
\Lambda 的敛散性. 
解.
ns
n
n
n
=-
+
-
+
+
-
=-
+
+
\Lambda 
,故lim
n
n
s
\to \infty 
=,级数收敛. 
\textcolor{lyred}{例}.讨论
(
)
n
n n
\infty 
=
+
\sum 
的收敛性. 
解.
(
)
1 1
n
n
n
n
k
k
k
s
k k
k
k
k
k
k
k
=
=
=








=
=
-
=
-
+
-
=








+
+
+
+
+








\sum 
\sum 
\sum 
1 1
2 2
n
n
k
k
k
k
k
k
n
n
=
=








-
+
-
=
-
+
-








+
+
+
+
+








\sum 
\sum 
,故 
lim
n
n
s
\to \infty 
=
,级数收敛. 
\textcolor{lyred}{例}.讨论
(
)
n
n
n
n
\infty 
=
+
-
++
\sum 
的敛散性. 
解.
(
)
n
n
n
k
k
s
k
k
k
k
k
k
k
=
=


=
+
-
++
=
-
=


+
+
+
++


\sum 
\sum 
n
n
-
+
+
+
+,故
lim
n
n
s
\to \infty 
=-
+,级数收敛. 
\textcolor{lyred}{例}.设
na \ge 
,讨论
(
)(
)
(
)
n
n
n
a
a
a
a
\infty 
=
+
+
+
\sum 
\Lambda 
的敛散性. 
解.由
(
)
(
)
(
)
(
)
(
)
(
)
1 1
n
n
n
n
n
a
u
a
a
a
a
a
a
-
+-
=
=
-
+
+
+
+
+
+
\Lambda 
\Lambda 
\Lambda 
,得 
(
)
(
)
n
n
s
a
a
=-
\le 
+
+
\Lambda 
,同时,\{\}
ns
单调增加,故级数收敛. 
\vspace{4pt}
{\bfseries \textcolor{lyred}{二.常数项级数的性质}}

\textcolor{lyred}{性质1}.设
n
n
u
\infty 
=\sum 
收敛,则
n
n
ku
\infty 
=\sum 
也收敛,且
n
n
n
n
ku
k
u
\infty 
\infty 
=
=
=
\sum 
\sum 
. 
\textcolor{lyred}{性质2}.设
n
n
u
\infty 
=\sum 
与
n
n
v
\infty 
=\sum 
均收敛,则
(
)
n
n
n
u
v
\infty 
=
\pm 
\sum 
也收敛,且 
(
)
n
n
n
n
n
n
n
u
v
u
v
\infty 
\infty 
\infty 
=
=
=
\pm 
=
\pm 
\sum 
\sum 
\sum 
. 
\textcolor{lyred}{注}.设
n
n
u
\infty 
=\sum 
收敛,
n
n
v
\infty 
=\sum 
发散,则
(
)
n
n
n
u
v
\infty 
=
\pm 
\sum 
发散. 
\textcolor{lyred}{例}.讨论
(
)
(
)
,
n
n
n
n
a
b
a b
a
b
\infty 
=
+
>
+
\sum 
的敛散性. 
解.由于
n
n
a
a
b
\infty 
=




+


\sum 
与
n
n
b
a
b
\infty 
=




+


\sum 
均收敛,故级数收敛. 
\textcolor{lyred}{性质3}.级数去掉,加上或改变有限多项,不改变其敛散性. 
\textcolor{lyred}{性质4}.收敛级数的项任意加括号后仍收敛,且和不变. 
\textcolor{lyred}{注}.加括号后发散的级数,本身一定发散,例如\textcolor{lyred}{调和级数}
n
n
\infty 
=\sum 
,按如下方式加括号 
n
n
v
v
-







+
+
+
+
+
+
+
+
+
+
+
+
+
+
=
+
+







+







\Lambda 
\Lambda 
\Lambda 
\Lambda 
\Lambda 
\Lambda , 
则
nv >
,发散,故原级数发散. 
\textcolor{lyred}{注}.加括号后收敛的级数,本身不一定收敛,例如(
)(
)
1 1
1 1
-
+
-
+\Lambda 收敛,而 
1 1 1 1
-+-+\Lambda 是发散的. 
\textcolor{lyred}{例}.设
(
)
n
n
n
u
u
\infty 
-
=
+
\sum 
收敛,且lim
n
n
u
\to \infty 
=
,证明:
n
n
u
\infty 
=\sum 
收敛. 
证.
(
)
n
n
k
k
k
s
u
u
-
=
=
+
\sum 
收敛,而
(
)
n
n
k
k
n
k
s
u
u
u
+
-
+
=
=
+
+
\sum 
,由lim
n
n
u
\to \infty 
=
,故 
lim
lim
n
n
n
n
s
s
+
\to \infty 
\to \infty 
=
也收敛,证毕. 
\vspace{4pt}
{\bfseries \textcolor{lyred}{三.收敛的必要条件}}

\textcolor{lyred}{定理}.设
n
n
u
\infty 
=\sum 
收敛,则lim
n
n
u
\to \infty 
=
. 
\textcolor{lyred}{注}.若lim
n
n
u
\to \infty 
\neq 
,则
n
n
u
\infty 
=\sum 
发散,例如
(
)
n
n
n
n
\infty 
=
-
+
\sum 
,
n
n
n
\infty 
=\sum 
,
(
)
n
n
n
n
n
\infty 
=
+
\sum 
. 
\textcolor{lyred}{注}.反之不对,例如\textcolor{lyred}{调和级数}
n
n
\infty 
=\sum 
发散,但是
lim
n
n
\to \infty 
=
. 
\textcolor{lyred}{补充练习 }
1.讨论
(
)
1 !
n
n
n
\infty 
=
+
\sum 
的收敛性. 
解.
(
)
(
)
(
)
1 1
1 !
!
1 !
1 !
n
n
n
u
s
n
n
n
n
+-
=
=
-
\Rightarrow 
=-
+
+
+
,故lim
n
n
s
\to \infty 
=,级数收敛. 
2.讨论
3 1
n
n
-
+
-
+
+
-
+
-
+
-
+
-
+
\Lambda 
\Lambda 的敛散性. 
解.加括号,得
n
n
n
n
n
n
n
\infty 
\infty 
\infty 
=
=
=


-
=
=


-
-
+


\sum 
\sum 
\sum 
,故级数发散. 
\vspace{6pt}
{\large \bfseries \textcolor{lyred}{第12.2 节 常数项级数的审敛法}}

\vspace{4pt}
{\bfseries \textcolor{lyred}{一.正项级数及其审敛法}}

\textcolor{lyred}{定义}.若
n
u \ge 
,则称
n
n
u
\infty 
=\sum 
为\textcolor{lyred}{正项级数}. 
{\bfseries \textcolor{lyred}{1.基本定理}}

\textcolor{lyred}{定理}.正项级数
n
n
u
\infty 
=\sum 
收敛\Leftrightarrow 部分和数列\{\}
ns
有界. 
\textcolor{lyred}{推论(柯西积分判别法)}.设
()
f x 是[
)
1,+\infty 上的单调减少连续函数,且
()
f x >
, 
令
()
nu
f n
=
,则
n
n
u
\infty 
=\sum 
与
()
f x dx
+\infty 
\int 
有相同的收敛性. 
\textcolor{lyred}{例}.讨论\textcolor{lyred}{p-级数}
p
p
p
p
n
n
n
\infty 
=
=+
+
+
+
+
\sum 
\Lambda 
\Lambda 的收敛性. 
解.考虑
p dx
x
+\infty 
\int 
的收敛性,当
p >时收敛,当
p \le 时发散. 
\textcolor{lyred}{例}.讨论\textcolor{lyred}{对数级数}
ln p
n
n
n
\infty 
=\sum 
的敛散性. 
解.因为
ln2
ln p
p
dx
du
x
x
u
+\infty 
+\infty 
=
\int 
\int 
,故当
p >时收敛,当
p \le 时发散. 
{\bfseries \textcolor{lyred}{2.比较审敛法}}

\textcolor{lyred}{定理}.设
n
n
u
\infty 
=\sum 
,
n
n
v
\infty 
=\sum 
为正项级数,(1)若
(
)
n
n
u
v
n
N
\le 
>
,
n
n
v
\infty 
=\sum 
收敛,则
n
n
u
\infty 
=\sum 
收敛; 
(2)若
(
)
n
n
u
v
n
N
\ge 
>
,
n
n
v
\infty 
=\sum 
发散,则
n
n
u
\infty 
=\sum 
发散. 
\textcolor{lyred}{推论}.设正项级数
n
n
u
\infty 
=\sum 
收敛,则
\alpha 
\forall 
>,
n
n
u\alpha 
\infty 
=\sum 
也收敛. 
\textcolor{lyred}{例}.讨论
ln 1
n
n
n
\infty 
=




-
+








\sum 
的收敛性. 
解.利用
(
)
(
)
ln 1
x
x
x x
x <
+
<
>
+
(见3.1),得
ln 1
n
n
n


<
+
<


+


,故 
(
)
ln 1
n
n
n
n
n n
n


<
-
+
<
-
=
<


+
+


,级数收敛. 
\textcolor{lyred}{例}.设
n
n
a
\infty 
=\sum 
与
n
n
b
\infty 
=\sum 
均收敛,证明
n
n
n
a b
\infty 
=\sum 
收敛. 
证.由
(
)
n
n
n
n
a b
a
b
\le 
\le 
+
,即得,证毕. 
\textcolor{lyred}{注}.特别地,若
n
n
a
\infty 
=\sum 
收敛,则
n
n
a
n
\infty 
=\sum 
收敛. 
\textcolor{lyred}{例}.讨论下列级数的收敛性 
(1)
sin
2n
n
n
\infty 
=\sum 
,(2)
(
)
ln
ln
n
n
n
\infty 
=\sum 
,(3)
ln !
n
n
\infty 
=\sum 
,(4)
ln
n
n
n\alpha 
\infty 
=\sum 
. 
解.(2)
(
)
ln
ln
lnln
lnln
ln
n
n
n
n
e
n
n
n
\cdots 
=
=
<
,收敛;(3)
ln !
ln
n
n
n
>
,发散; 
(4)当
\alpha \le 时, ln
n
n
n
\alpha 
\alpha 
>
,发散;当
\alpha >时,取
\varepsilon >
,使得
\alpha 
\varepsilon 
=+
,则 
(
)
1 2
ln
ln
ln
n
n
n
n
N
n
n
n
n
n
\alpha 
\varepsilon 
\varepsilon 
\varepsilon 
\varepsilon 
+
+
+
=
=
\cdots 
\le 
>
,收敛. 
\textcolor{lyred}{例}.设
n
n
n
v
u
w
\le 
\le 
,且
n
n
v
\infty 
=\sum 
与
n
n
w
\infty 
=\sum 
均收敛,证明:
n
n
u
\infty 
=\sum 
收敛. 
证.0
n
n
n
n
u
v
w
v
\le 
-
\le 
-
,
(
)
n
n
n
w
v
\infty 
=
-
\sum 
收敛
(
)
n
n
n
u
v
\infty 
=
\Rightarrow 
-
\sum 
收敛,即得,证毕. 
\textcolor{lyred}{推论}.设
n
n
u
\infty 
=\sum 
与
n
n
v
\infty 
=\sum 
均为正项级数. 
(1)若存在
k >
,使得
(
)
n
n
u
kv
n
N
\le 
>
,
n
n
v
\infty 
=\sum 
收敛,则
n
n
u
\infty 
=\sum 
收敛; 
(2)若存在
l >
,使得
(
)
n
n
u
lv
n
N
\ge 
>
,
n
n
v
\infty 
=\sum 
发散,则
n
n
u
\infty 
=\sum 
也发散; 
(3)若存在
k >
,
l >
,使得
(
)
n
n
n
kv
u
lv
n
N
\ge 
\ge 
>
,则
n
n
u
\infty 
=\sum 
与
n
n
v
\infty 
=\sum 
的收敛性相同. 
{\bfseries \textcolor{lyred}{3.比较法的极限形式}}

\textcolor{lyred}{定理}.设
n
n
u
\infty 
=\sum 
与
n
n
v
\infty 
=\sum 
为正项级数,lim
n
n
n
u
l
v
\to \infty 
=. 
(1)当0
l
<<+\infty 时,
n
n
u
\infty 
=\sum 
与
n
n
v
\infty 
=\sum 
的收敛性相同; 
(2)当
l =
时,
n
n
v
\infty 
=\sum 
收敛,则
n
n
u
\infty 
=\sum 
收敛; 
(3)当l =+\infty 时,
n
n
v
\infty 
=\sum 
发散,则
n
n
u
\infty 
=\sum 
发散. 
\textcolor{lyred}{推论}.设
n
n
u
\infty 
=\sum 
,
n
n
v
\infty 
=\sum 
为正项级数,
(
)
n
n
u
v
n \to \infty 
:
,则
n
n
u
\infty 
=\sum 
与
n
n
v
\infty 
=\sum 
的收敛性相同. 
\textcolor{lyred}{例}.讨论
ln 1
n
n
n
\infty 
=




-
+








\sum 
的收敛性. 
解.由于
(
)
(
)
ln 1
x
x
x
o x
x
-
+
=
+
:
,故
ln 1
n
n
n
\infty 
=




-
+








\sum 
收敛. 
\textcolor{lyred}{例}.讨论
(
)
1 cos
p
n
n
n
n
\infty 
=
+
-
-
\sum 
的收敛性. 
解.
1 2
p
n
p
u
n
n
n
+

\cdots 
\cdots 
=
\cdots 




:
,故当
p >
时收敛,当
p \le 
时发散. 
\textcolor{lyred}{例}.讨论下列级数的收敛性 
(1)
n
n
n
\infty 
=
+
-
\sum 
,(2)
ln
n
n
n
n
n
n
\infty 
=
+
+
-
-
\sum 
,(3)
(
)
1 sin
n
n
n
n
\infty 
=
-
+
\sum 
, 
(4)
ln 1
n
n
n
\infty 
=


+


+


\sum 
,(5)
n
n
n
\infty 
+
=


-






\sum 
. 
解.(5)
ln
ln
ln
n
n
n
n
n
n
e
n
n
+
+
-=
-
+
:
:
,收敛. 
\textcolor{lyred}{例}.设
na \ge 
,且
2 sin
lim
n
n
n
n
n
a
\to \infty 


\cdots 
=




,证明
n
n
a
\infty 
=\sum 
收敛. 
证.当n \to \infty 时,
2 sin
n
n
n
a
n
:
,而
lim 2 sin
n
n
n
\to \infty 
=
,故当n
N
>
时,
2 sin
n
n
n
n
<
,于是 
2 sin
n
n
n
n
\infty 
=\sum 
收敛,证毕. 
\textcolor{lyred}{推论(极限审敛法)}.设lim
p
n
n
n u
l
\to \infty 
=,(1)若0
l
\le <+\infty ,且
p >,则
n
n
u
\infty 
=\sum 
收敛; 
(2)若0
l
<\le +\infty ,且
p \le ,则
n
n
u
\infty 
=\sum 
发散. 
{\bfseries \textcolor{lyred}{4.比值审敛法(达朗贝尔判别法)}}

\textcolor{lyred}{定理}.设
n
n
u
\infty 
=\sum 
为正项级数,
lim
n
n
n
u
u
\rho 
+
\to \infty 
=
,则(1)当
\rho <时,级数收敛; 
(2)当
\rho >(包括\rho =+\infty )时,级数发散,此时lim
n
n
u
\to \infty 
=\infty . 
\textcolor{lyred}{例}.讨论下列级数的收敛性 
(1)
(
)
k
n
n
n
a
a
\infty 
=
>
\sum 
,(2)
(
)
!
n
n
a
a
n
\infty 
=
>
\sum 
,(3)
!
n
n
n
n
\infty 
=\sum 
,(4)
1 3
n
n
n
n
\infty 
=
-
\sum 
. 
解.(1)
lim
lim
k
n
n
n
n
u
n
u
a
n
a
+
\to \infty 
\to \infty 
+


=
=




,收敛;(2)
lim
lim
n
n
n
n
u
a
u
n
+
\to \infty 
\to \infty 
=
=
+
,收敛; 
(3)
lim
lim
n
n
n
n
n
u
n
u
n
e
+
\to \infty 
\to \infty 


=
=


+


,收敛;(4)
lim
lim 3
n
n
n
n
n
n
n
n
u
u
+
+
+
\to \infty 
\to \infty 
-
=
=
-
,收敛. 
\textcolor{lyred}{注}.lim
k
n
n
n
a
\to \infty 
=
,lim
!
n
n
a
n
\to \infty 
=
,
!
lim
n
n
n
n
\to \infty 
=
,即n \to \infty 时,
!
k
n
n
n
a
n
n
=
=
=
. 
\textcolor{lyred}{例}.讨论
!
n
n
n
e
n
n
\infty 
=
\cdots 
\sum 
的收敛性. 
解.
lim
lim
n
n
n
n
n
u
e
u
n
+
\to \infty 
\to \infty 
=
=


+




,无法判断,但是考虑到
n
e
n


+
<




,即
n
n
u
u
+>\Rightarrow  
lim
n
n
u
\to \infty 
\neq 
,故发散. 
{\bfseries \textcolor{lyred}{5.根值审敛法(柯西判别法)}}

\textcolor{lyred}{定理}.设
n
n
u
\infty 
=\sum 
为正项级数,lim n
n
n
u
\rho 
\to \infty 
=
,则(1)当
\rho <时,级数收敛; 
(2)当
\rho >(包括\rho =+\infty )时,级数发散,此时lim
n
n
u
\to \infty 
=\infty . 
\textcolor{lyred}{例}.讨论下列级数的收敛性 
(1)
n
n
n
n
\infty 
=
+




-


\sum 
,(2) 
lnn
n
n
n
\infty 
=\sum 
,(3)
n
n
n
n
\infty 
=


+




\sum 
,(4)
1 3
n
n
n
n
\infty 
=
-
\sum 
. 
解.(1)
lim
n
n
n
u
\to \infty 
=
,收敛;(2) lim n
n
n
u
\to \infty 
=+\infty ,发散; 
(3)lim
n
n
n
e
u
\to \infty 
=
<,收敛;(4)
lim
lim
n
n
n
n
n
n
n
u
\to \infty 
\to \infty 
=
=
-
,收敛. 
\vspace{4pt}
{\bfseries \textcolor{lyred}{二.交错级数及其审敛法}}

\textcolor{lyred}{交错级数}:正负相间的级数
(
)
n
n
n
u
\infty 
-
=
-
\sum 
,其中
n
u >
. 
\textcolor{lyred}{定理(莱布尼茨定理)}.设
n
u 单调减少,且lim
n
n
u
\to \infty 
=
,则交错级数
(
)
n
n
n
u
\infty 
-
=
-
\sum 
收敛, 
且和
s
u
\le 
,余项nr 的绝对值
n
n
r
u +
\le 
. 
\textcolor{lyred}{例}.当
p >
时,
(
)
n
p
p
p
p
n
n
\infty 
-
=
-
=-
+
-
+
\sum 
\Lambda 收敛. 
\textcolor{lyred}{例}.讨论下列级数的收敛性 
(1)
(
)
ln
n
n
n
n
\infty 
=
-
\sum 
,(2)
(
)
sin
n
n
a
\pi 
\infty 
=
+
\sum 
. 
解.(1)
(
)
ln
ln
x
x
x
e
x
x x
'
-

=
<
>




,又
ln
lim
x
x
x
\to +\infty 
=
,收敛; 
(2)
(
)
(
)
(
)
(
)
sin
sin
sin
n
n
a
n
a
n
a
n
n
a
n
\pi 
\pi 
\pi 
\pi 
+
=-
+
-
=-
+
+
,收敛. 
\textcolor{lyred}{例}.设
n
u 为单调减少正数列,
(
)
n
n
n
u
\infty 
=
-
\sum 
发散,证明
n
n
n
u
\infty 
=




+


\sum 
收敛. 
证.
n
u 单调有界,故收敛,设lim
n
n
u
a
\to \infty 
=
\ge 
,由于
(
)
n
n
n
u
\infty 
=
-
\sum 
发散,故
a >
,于是 
由根值法,或者
n
n
n
u
a




\le 
\le 




+
+




,即得,证毕. 
\vspace{4pt}
{\bfseries \textcolor{lyred}{三.绝对收敛与条件收敛}}

\textcolor{lyred}{定义}.若
n
n
u
\infty 
=\sum 
收敛,则称
n
n
u
\infty 
=\sum 
\textcolor{lyred}{绝对收敛},例如
(
)
(
)
n
p
n
p
n
\infty 
-
=
-
>
\sum 
; 
若
n
n
u
\infty 
=\sum 
发散,
n
n
u
\infty 
=\sum 
收敛,则称
n
n
u
\infty 
=\sum 
\textcolor{lyred}{条件收敛},例如
(
)
n
n
n
\infty 
-
=
-
\sum 
. 
\textcolor{lyred}{例}.讨论
(
)
(
)
sin
n
n
n\alpha \alpha 
\infty 
-
=
-
>
\sum 
的收敛性. 
解.当
\alpha >时,由
sin n
n
\alpha 
\alpha 
:
,级数绝对收敛;当
\alpha \le 时,级数条件收敛. 
\textcolor{lyred}{注}.对于
(
)
(
)
ln 1
n
n
n\alpha 
\alpha 
\infty 
-
=


-
+
>




\sum 
,有类似的结论. 
\textcolor{lyred}{定理}.设
n
n
u
v
\le 
,
n
n
v
\infty 
=\sum 
收敛,则
n
n
u
\infty 
=\sum 
绝对收敛. 
\textcolor{lyred}{例}.级数
sin
n
n
n
\alpha 
\infty 
=\sum 
绝对收敛.\textcolor{lyred}{ }
\textcolor{lyred}{例}.设
n
n
u
\infty 
=\sum 
与
n
n
v
\infty 
=\sum 
绝对收敛,证明
(
)
n
n
n
u
v
\infty 
=
\pm 
\sum 
绝对收敛. 
证.
n
n
n
n
u
v
u
v
\pm 
\le 
+
,由比较法即得,证毕. 
\textcolor{lyred}{例}.讨论级数
(
)
n
n
n
n
\infty 
=


-
+






\sum 
的收敛性. 
解.
n
n
\infty 
=\sum 
绝对收敛,
(
)
n
n
n
\infty 
=
-
\sum 
条件收敛,故级数条件收敛. 
\textcolor{lyred}{定理}.绝对收敛的级数一定收敛. 
\textcolor{lyred}{证}.令
,
0,
n
n
n
n
n
n
u
u
u
u
u
u
+
\ge 
+

=
=
<

,则0
n
n
u
u
+
\le 
\le 
,故
n
n
u
\infty 
+
=\sum 
收敛,而
n
n
n
u
u
u
+
=
-
, 
故
n
n
u
\infty 
=\sum 
收敛,证毕. 
\textcolor{lyred}{注}.又令
0,
,
n
n
n
n
n
n
u
u
u
u
u
u
-
\ge 
-

=
=-
<

,则
n
n
n
u
u
u
+
-
=
-
,
n
n
n
u
u
u
+
-
=
+
,故 
n
n
u
\infty 
=\sum 
绝对收敛
n
n
u
\infty 
+
=
\Leftrightarrow \sum 
及
n
n
u
\infty 
-
=\sum 
均收敛;
n
n
u
\infty 
=\sum 
条件收敛
n
n
u
\infty 
+
=
\Rightarrow \sum 
及
n
n
u
\infty 
-
=\sum 
均发散. 
\textcolor{lyred}{注}.一般地,
n
n
u
\infty 
=\sum 
发散不能推出
n
n
u
\infty 
=\sum 
发散,但是,若
n
n
u
\infty 
=\sum 
发散是由于比值法或 
根值法中的
\rho >,则
n
n
u
\infty 
=\sum 
也发散,因为此时lim
n
n
u
\to \infty 
=\infty . 
\textcolor{lyred}{例}.讨论下列级数的收敛性 
(1)
(
)
(
)
n n
n
n
n
\infty 
+
=
-
\sum 
,(2)
(
)
!
n
n
n
n
n
n
-
\infty 
=
-
\sum 
,(3)
(
)
n
n
n
n
n
\infty 
=
-


+




\sum 
, 
(4)
(
)(
)
0,
n
s
n
s
n
\alpha 
\alpha 
\infty 
=
-
>
>
\sum 
. 
解.(1)
lim
n
n
n
u
u
+
\to \infty 
=
,或
lim
n
n
n
u
\to \infty 
=
<,绝对收敛; 
(2)
(
)
(
)
lim
lim3
n
n
n
n
n
n
u
n
n
u
e
n
+
+
\to \infty 
\to \infty 
=
+
=
>
+
,发散; 
(3)
lim
lim
n
n
n
n
n
e
u
n
\to \infty 
\to \infty 


=
+
=
>




,发散; 
(4) lim n
n
n
u
\alpha 
\to \infty 
=
,故0
\alpha 
<
<时绝对收敛,
\alpha >时发散,而当
\alpha =时, 
若0
s
<
\le ,则条件收敛,若
s >,则绝对收敛. 
\vspace{4pt}
{\bfseries \textcolor{lyred}{四.绝对收敛级数的性质}}

{\bfseries \textcolor{lyred}{1.交换律}}

\textcolor{lyred}{定理}.设
n
n
u
\infty 
=\sum 
绝对收敛,则任意交换各项顺序所得到的新级数
*
n
n
u
\infty 
=\sum 
仍绝对收敛, 
且和不变. 
\textcolor{lyred}{注}.对于条件收敛的级数,交换律不成立. 
\textcolor{lyred}{例}.
(
)
(
)
n
n
n
s
n
n
\infty 
-
-
=
-
=-
+
-
+
+-
+
=
\sum 
\Lambda 
\Lambda 
,重排如下: 
k
k
k
-
-
+
-
-
+
-
-
+
-
+
-
-
+
-
-
\Lambda 
\Lambda ,则 
*
n
n
n
n
k
k
k
s
k
k
k
k
k
k
k
=
=
=






=
-
-
=
-
=
-
=






-
-
-
-






\sum 
\sum 
\sum 
n
s
n
n


-
+
-
+
-
=


-


\Lambda 
,故
*
lim
n
n
s
s
\to \infty 
=
,而
*
*
n
n
s
s
n
-=
+
, 
*
*
n
n
s
s
n
n
-=
+
+
-
,故
*
*
lim
lim
n
n
n
n
s
s
s
-
-
\to \infty 
\to \infty 
=
=
,级数收敛于1
2 s ; 
再重排如下: 
k
k
k
k
-
-
-
+
-
-
-
+
+
-
-
-
+
-
-
-
\Lambda 
\Lambda ,则 
*
n
n
n
k
k
s
k
k
k
k
k
k
k
=
=




=
-
-
-
=
-
-
=




-
-
-
-
-




\sum 
\sum 
n
n
n
k
k
s
k
k
k
k
=
=




-
=
-
=




-
-




\sum 
\sum 
,故
*
lim
n
n
s
s
\to \infty 
=
,而 
*
n
n
s
s
n
-=
+
,
*
n
n
s
s
n
n
-=
+
+
-
,
*
n
n
s
s
n
n
n
-=
+
+
+
-
-
, 
故
*
*
*
lim
lim
lim
n
n
n
n
n
n
s
s
s
s
-
-
-
\to \infty 
\to \infty 
\to \infty 
=
=
=
,级数收敛于1
4 s . 
{\bfseries \textcolor{lyred}{2.级数的乘法}}

\textcolor{lyred}{正方形法}:
(
)(
)
1 1
1 2
2 2
2 1
1 3
2 3
3 3
2 3
1 3
u v
u v
u v
u v
u v
u v
u v
u v
u v
+
+
+
+
+
+
+
+
+\Lambda . 
\textcolor{lyred}{对角线法(柯西乘积)}:
(
)
(
)
1 1
1 2
2 1
n
n
u v
u v
u v
u v
u v
+
+
+
+
+
+
+
\Lambda 
\Lambda 
\Lambda . 
\textcolor{lyred}{定理}.设
n
n
u
\infty 
=\sum 
,
n
n
v
\infty 
=\sum 
绝对收敛,则其柯西乘积也绝对收敛,和为
n
n
n
n
s
u
v
\infty 
\infty 
=
=
=
\cdots 
\sum 
\sum 
. 
\textcolor{lyred}{补充练习 }
1.设数列
n
na 收敛,证明
n
n
a
\infty 
=\sum 
收敛. 
证.
n
na 有界,设
n
na
M
\le 
,则
n
M
a
n
\le 
,故
n
n
a
\infty 
=\sum 
收敛,证毕. 
2.讨论
sin 3n
n
n
\infty 
=\sum 
的收敛性. 
解.
n
n
n
u :
,再由比值法,级数收敛. 
3.设
(
)
n
n
n
n
a
\infty 
=
-
\sum 
收敛,讨论\sum 
\infty 
=1
n
n
a 的敛散性. 
解.
(
)
lim
n
n
n
n
a
\to \infty -
=
,故
n
n
n
n
M
a
M
a
\le 
\Rightarrow 
\le 
,故绝对收敛. 
4.设
()
[
]
0,2
f x
C
\pi 
\in 
,记
()
cos
na
f x
nxdx
\pi 
=\int 
,
1,2,
n =
\Lambda ,证明: 
(1)若
()
f x 在[
]
0,2\pi 上有连续的一阶导数,则
n
n
a
\infty 
=\sum 
收敛; 
(2)若
()
f x 在[
]
0,2\pi 上有连续的二阶导数,则
n
n
a
\infty 
=\sum 
绝对收敛. 
证.(1)
()(
)
()
sin
sin
na
f x d
nx
f
x
nxdx
n
n
\pi 
\pi 
'
=
=-
\int 
\int 
,故
n
M
a
n
\le 
,即得; 
(2)
()(
)
(
)
()
()
cos
cos
n
f
f
a
f
x d
nx
f
x
nxdx
n
n
n
\pi 
\pi 
\pi 
'
'
-
'
''
=
=
-
\int 
\int 
,故
n
M
a
n
\le 
, 
即得,证毕. 
5.设
(
)
na
a
nd d
=
+
>
为等差数列,证明
n
na
\infty 
=\sum 
发散. 
证.
1 1
na
a
nd
d n
=
+
:
,证毕. 
6.设
()
f x 在
x =
处二阶可导,
()
lim
x
f x
x
\to 
=
,证明
n
f
n
\infty 
=






\sum 
绝对收敛. 
证.
()
()
()
()
(
)
()
(
)
f
f
f x
f
f
x
x
o x
x
o x
''
''
'
=
+
+
+
=
+
,于是 
()
lim
n
f
n
f
n
\to \infty 
''

=




,故
n
f
n
+\infty 
=






\sum 
收敛,证毕. 
7.设
n
u \neq 
,且lim
n
n
n
u
\to \infty 
=,讨论
(
)
n
n
n
n
u
u
\infty 
-
=
+


-
+




\sum 
的敛散性. 
解.
lim
lim
n
n
n
n
n
u
u
n
\to \infty 
\to \infty 


=
\cdots 
=




,
(
)
n
k
n
k
k
k
n
s
u
u
u
u
u
+
=
+
+


=
-
+
=
\pm 
\to 




\sum 
,故 
级数收敛,而
lim
n
n
n
n
n
n u
u
u
u
n
\to \infty 
+
+


+
=
\Rightarrow 
+




:
,故级数条件收敛. 
8.讨论
(
)
(
)
(
)
1 3 5
2 4 6
n
n
n
n
\infty 
-
=
\cdots \cdots 
-
-
\cdots \cdots 
\sum 
\Lambda 
\Lambda 
的收敛性. 
解.设
(
)
(
)
1 3 5
2 4 6
n
n
u
n
\cdots \cdots 
-
=
\cdots \cdots 
\Lambda 
\Lambda 
,则
n
n
n
n
u
u
u
n
+
+
=
<
+
,单调减少,又利用 
(
)(
)
(
)(
)
n
n
n
n
n
-
+
+
=
>
-
+
,得到 
(
)
1 3 5
n
n
u
n
n
n
\cdots \cdots 
-
<
=
\to 
\cdots 
\cdots 
\cdots 
\cdots 
\cdots 
-\cdots 
+
+
\Lambda 
\Lambda 
,级数收敛; 
(
)
(
)
3 5 7
3 5 7
2 4 6
2 4 6
n
n
n
u
n
n
n
n
n
\cdots \cdots 
-
+
=
=
\cdots 
\cdots 
\cdots 
>
\cdots \cdots 
-
\cdots 
+
+
\Lambda 
\Lambda 
\Lambda 
,故
n
n
u
\infty 
=\sum 
发散,即 
(
)
n
n
n
u
\infty 
-
=
-
\sum 
条件收敛. 
\vspace{6pt}
{\large \bfseries \textcolor{lyred}{第12.3 节 幂级数}}

\vspace{4pt}
{\bfseries \textcolor{lyred}{一.函数项级数}}

设
()
\{
\}
nu
x
为I 上函数序列,则表达式
()
n
n
u
x
\infty 
=\sum 
称为I 上的\textcolor{lyred}{函数项(无穷)级数}; 
若
(
)
n
n
u
x
\infty 
=\sum 
收敛,则称
0x 为级数的\textcolor{lyred}{收敛点},反之称为\textcolor{lyred}{发散点},收敛点全体构成的 
集合称为\textcolor{lyred}{收敛域},发散点全体构成的集合称为\textcolor{lyred}{发散域},在收敛域内 
()
()
n
n
s x
u
x
\infty 
=
=\sum 
称为该函数项级数的\textcolor{lyred}{和函数},
()
()
()
n
n
r
x
s x
s
x
=
-
为\textcolor{lyred}{余项}. 
\textcolor{lyred}{例}.
n
n
n
x
x
n
-
\infty 
=
+
\sum 
的收敛域为\{
\}
x =\pm 
. 
\vspace{4pt}
{\bfseries \textcolor{lyred}{二.幂级数及其收敛性}}

\textcolor{lyred}{幂级数}
(
)
(
)
(
)
n
n
a
a
x
x
a
x
x
a
x
x
+
-
+
-
+
+
-
+
\Lambda 
\Lambda 是最常见的函数项级数, 
我们只讨论特殊形式
n
n
n
n
n
a x
a
a x
a x
a x
\infty 
=
=
+
+
+
+
+
\sum 
\Lambda 
\Lambda ,其中
na 称为\textcolor{lyred}{系数}. 
\textcolor{lyred}{注}.幂级数
n
n
n
a x
\infty 
=\sum 
在
x =
处必收敛,故收敛域一定非空. 
{\bfseries \textcolor{lyred}{1.幂级数收敛域的结构}}

\textcolor{lyred}{例}.几何级数
n
n
x
\infty 
=\sum 
,
()
n
n
k
n
k
x
s
x
x
x
-
=
-
=
=
-
\sum 
,当
x <时收敛, ()
s x
x
=-
, 
当
x \ge 时发散,故其收敛域为(
)
1,1
-
. 
\textcolor{lyred}{定理(Abel 定理)}.(1)设
n
n
n
a x
\infty 
=\sum 
收敛,则当
x
x
<
时,
n
n
n
a x
\infty 
=\sum 
绝对收敛; 
(2)设
n
n
n
a x
\infty 
=\sum 
发散,则当
x
x
>
时,
n
n
n
a x
\infty 
=\sum 
发散. 
\textcolor{lyred}{推论}.设A 为
n
n
n
a x
\infty 
=\sum 
的收敛域,令
\{
\}
sup
:
R
x
x
A
=
\in 
(可以为\infty ),则 
(1)当x
R
<
时,
n
n
n
a x
\infty 
=\sum 
绝对收敛;(2)当x
R
>
时,
n
n
n
a x
\infty 
=\sum 
发散. 
\textcolor{lyred}{定义}.称R 为
n
n
n
a x
\infty 
=\sum 
的\textcolor{lyred}{收敛半径},(
)
,
R R
-
为\textcolor{lyred}{收敛区间}. 
\textcolor{lyred}{例}.几何级数
n
n
x
\infty 
=\sum 
的收敛半径
R =,收敛区间为(
)
1,1
-
. 
\textcolor{lyred}{注}.当
R =
时只在
x =
处收敛, R =+\infty 时在(
)
,
-\infty +\infty 上处处收敛. 
\textcolor{lyred}{注}.一般地,对于幂级数
(
)
n
n
n
a
x
x
\infty 
=
-
\sum 
,也有收敛半径R 的概念,当
x
x
R
-
<
时, 
级数绝对收敛,当
x
x
R
-
>
时,级数发散. 
{\bfseries \textcolor{lyred}{2.收敛半径的求法}}

\textcolor{lyred}{定理(系数模比值法)}.设有
n
n
n
a x
\infty 
=\sum 
,若
na \neq 
,且
lim
n
n
n
a
a
\rho 
+
\to \infty 
=
,则
R
\rho 
=
. 
\textcolor{lyred}{例}.求下列幂级数的收敛半径 
(1)
(
)
n
n
n
x
n
\infty 
-
=
-
\sum 
,(2)
!
n
n
x
n
\infty 
=\sum 
,(3)
!
n
n
n x
\infty 
=\sum 
,(4)
(
)
n
n
n
x
n
\infty 
=
-
\sum 
. 
解.(1)
lim
lim
n
n
n
n
a
n
R
a
n
\rho 
+
\to \infty 
\to \infty 
=
=
=\Rightarrow 
=
+
; 
(2)
lim
lim
n
n
n
n
a
R
a
n
\rho 
+
\to \infty 
\to \infty 
=
=
=
\Rightarrow 
=+\infty 
+
; 
(3)
(
)
lim
lim
n
n
n
n
a
n
R
a
\rho 
+
\to \infty 
\to \infty 
=
=
+
=\infty \Rightarrow 
=
; 
(4)考虑
n
n
n
t
n
\infty 
=\sum 
,
lim
lim
n
n
n
n
a
n
R
a
n
\rho 
+
\to \infty 
\to \infty 
=
=
=
\Rightarrow 
=
+
. 
\textcolor{lyred}{例}.求
n
n
n
n
x
\infty 
=
+
\sum 
的收敛区间. 
解.令
t
x
=
,则级数变为
n
n
n
n
t
\infty 
=
+
\sum 
,
lim
n
n
n
a
R
a
\rho 
+
\to \infty 
=
=
\Rightarrow 
=
,故收敛区间为 
t
-<<
,即
x
-
<
<
. 
\textcolor{lyred}{注}.一般地,对于
()
n
n
u
x
\infty 
=\sum 
,可以由
()
()
lim
n
x
n
u
x
u
x
+
\to \infty 
<\Rightarrow 收敛区间. 
\textcolor{lyred}{例}.求
(
)
(
)
(
)
!
!
n
n
n
x
n
\infty 
=
+
\sum 
的收敛区间. 
解.
()
()
(
)
(
)
(
)
(
)(
)
(
)(
)
(
)
(
)
2 !
!
lim
lim
lim
!
1 !
n
x
x
x
n
u
x
n
n
n
n
x
x
u
x
n
n
n
+
\to \infty 
\to \infty 
\to \infty 
+
+
+
=
\cdots 
\cdots 
+
=
\cdots 
+
=
+
+




(
)
x +
,故(
)
x
x
x
+
<\Rightarrow -
<
+<
\Rightarrow -
<
<-
. 
\textcolor{lyred}{例}.求
(
)
n
n
n
n
x
n
\infty 
=


+-


\sum 
的收敛区间. 
解.考虑
n
n
n
x
n
-
\infty 
-
=
-
\sum 
,
n
n
n
x
n
\infty 
=\sum 
,前者的
R =
,后者的
R =
,故当
x <
时, 
原级数
n
n
n
n
n
n
x
x
n
n
-
\infty 
\infty 
-
=
=
=
+
-
\sum 
\sum 
绝对收敛,而当
x >
时,
lim 2
n
n
n
x
n
\to \infty 
=\infty ,故 
原级数发散,因此收敛区间为
1 1
,
4 4


-




. 
\textcolor{lyred}{定理(系数模根值法)}.设有
n
n
n
a x
\infty 
=\sum 
,若lim n
n
n
a
\rho 
\to \infty 
=
,则
R
\rho 
=
. 
\textcolor{lyred}{例}.求
n
n
n
x
n
\infty 
=


+




\sum 
的收敛区间. 
解.
lim
lim 1
n
n
n
n
n
a
e
R
n
e
\rho 
\to \infty 
\to \infty 


=
=
+
=
\Rightarrow 
=




,故收敛区间为
1 1
,
e e


-




. 
\textcolor{lyred}{例}.求
(
)
n
n
n
x
n
\infty 
+
=\sum 
的收敛区间. 
解.令
t
x
=
,考虑级数
4n
n
n
n
t
n
\infty 
=\sum 
,则
lim
n
n
n
a
R
\rho 
\to \infty 
=
=
\Rightarrow 
=+\infty ,故 
(
)
(
)
,
,
t
x
\in -\infty +\infty \Rightarrow 
\in -\infty +\infty . 
\vspace{4pt}
{\bfseries \textcolor{lyred}{三.幂级数的运算}}

\textcolor{lyred}{定理}.设
n
n
n
a x
\infty 
=\sum 
,
n
n
n
b x
\infty 
=\sum 
的收敛区间分别为(
)
,
R R
-
,(
)
,
R R
-
,若 
\{
\}
min
,
x
R
R R
<
=
,则(1)
(
)
n
n
n
n
n
n
n
n
n
n
a x
b x
a
b
x
\infty 
\infty 
\infty 
=
=
=
\pm 
=
\pm 
\sum 
\sum 
\sum 
; 
(2)
n
n
n
n
n
n
n
n
n
a x
b x
c x
\infty 
\infty 
\infty 
=
=
=
\cdots 
=
\sum 
\sum 
\sum 
,其中
n
n
k
n k
k
c
a b -
=
=\sum 
. 
\textcolor{lyred}{定理}.设幂级数
n
n
n
a x
\infty 
=\sum 
的和函数为()
s x ,则\textcolor{lyred}{ }
{\bfseries \textcolor{lyred}{(1)(连续性). }}

s x 在
n
n
n
a x
\infty 
=\sum 
的收敛域内连续,即
lim
n
n
n
n
x
x
n
n
a x
a x
\infty 
\infty 
\to 
=
=
=
\sum 
\sum 
; 
{\bfseries \textcolor{lyred}{(2)(逐项积分). }}

s x 在
n
n
n
a x
\infty 
=\sum 
的收敛域内可积,且
()
x
x
n
n
n
s x dx
a x dx
\infty 
=
=\sum 
\int 
\int 
; 
{\bfseries \textcolor{lyred}{(3)(逐项求导). }}

s x 在
n
n
n
a x
\infty 
=\sum 
的收敛区间内可导,且()
(
)
n
n
n
s x
a x
\infty 
=
'
'
=\sum 
. 
\textcolor{lyred}{注}.幂级数通过逐项求导,逐项积分后得到的新级数收敛半径不变,但在收敛区间 
端点处的敛散性可能会有改变,例如: 
(
)
n
n
n
x
n
-
\infty 
=
-
\sum 
的收敛域为(
]
1,1
-
,而
(
)
n
n
n
x
\infty 
-
-
=
-
\sum 
的收敛域为(
)
1,1
-
. 
\textcolor{lyred}{例}.求
(
)
n
n
n
x
\infty 
=
-
\sum 
的和函数. 
解.
lim
n
n
n
a
R
a
\rho 
+
\to \infty 
=
=
\Rightarrow 
=
,收敛域为
1 1
,
2 2


-




,则在
1 1
,
2 2


-




上 
()
(
)
(
)(
)
1 2
1 2
n
n
n
n
n
n
n
n
x
x
x
s x
x
x
x
x
x
x
x
\infty 
\infty 
\infty 
=
=
=
=
-
=
-
=
-
=
-
-
-
-
\sum 
\sum 
\sum 
. 
\textcolor{lyred}{例}.求
n
n
nx
\infty 
=\sum 
的和函数. 
解.
lim
n
n
n
a
R
\rho 
\to \infty 
=
=\Rightarrow 
=,收敛域为(
)
1,1
-
,设()
n
n
s x
nx
\infty 
=
=\sum 
,则在(
)
1,1
-
上, 
()
(
)
n
n
n
n
x
x
s x
x
nx
x
x
x
x
x
\infty 
\infty 
-
=
=
'
'




=
=
=
=




-


-


\sum 
\sum 
. 
\textcolor{lyred}{例}.求
n
n
n x
\infty 
=\sum 
的和函数. 
解.
lim
n
n
n
a
R
a
\rho 
+
\to \infty 
=
=\Rightarrow 
=,收敛域为(
)
1,1
-
,设()
n
n
s x
n x
\infty 
=
=\sum 
,则在(
)
1,1
-
上, 
()
(
)
n
n
n
n
n
n
n
n
n
n
s x
x
n x
x
n n
x
x
nx
x
x
x
x
\infty 
\infty 
\infty 
\infty 
\infty 
-
-
-
+
=
=
=
=
=
''
'




=
=
+
-
=
-
=








\sum 
\sum 
\sum 
\sum 
\sum 
(
)
(
)
(
)
x
x
x
x
x
x
x
x
x
x
x
x
x
x
x
x
x
''
'
''
'


+






-
=
-
=
-
=








-
-
-
-






-
-
-


. 
\textcolor{lyred}{例}.求
n
n
x
x
x
n
\infty 
=
=+
+
+
+
\sum 
\Lambda 的和函数. 
解.
lim
n
n
n
a
R
a
\rho 
+
\to \infty 
=
=\Rightarrow 
=,收敛域为[
)
1,1
-
,设()
n
n
x
s x
n
\infty 
=
=
+
\sum 
,则在(
)
1,1
-
上, 
()
()
n
n
n
n
x
xs x
xs x
x
n
x
+
\infty 
\infty 
=
=
'
=
\Rightarrow 
=
=




+
-
\sum 
\sum 
,故
()
()
()
x
xs x
s
xs x
dx
'
-
=
=




\int 
(
)
()
(
)(
)
ln 1
ln 1
x
x
dx
x
s x
x
x
x
-
=-
-
\Rightarrow 
=-
\neq 
-
\int 
,而()
s
=. 
\textcolor{lyred}{例}.求
(
)
(
)
n
n
n
x
n n
+
\infty 
+
=
-
+
\sum 
的和函数.  
解.
lim
n
n
n
a
R
a
\rho 
+
\to \infty 
=
=\Rightarrow 
=,收敛域为[
]
1,1
-
,设()
(
)
(
)
n
n
n
x
s x
n n
+
\infty 
+
=
=
-
+
\sum 
,则 
在(
)
1,1
-
上, ()
(
)
n
n
n
x
s
x
n
\infty 
+
=
'
=
-
\sum 
,
()
(
)
(
)
n
n
n
n
n
s
x
x
x
x
\infty 
\infty 
+
-
-
=
=
''
=
-
=
-
=+
\sum 
\sum 
, 
故()
()
(
)
ln 1
x
s
x
s
dx
x
x
'
'
-
=
=
+
+
\int 
, ()
()
(
)
ln 1
x
s x
s
x dx
-
=
+
=
\int 
(
)
(
)
ln 1
x
x
x
+
+
-
,
(
]
1,1
x\in -
,而(
)
()
lim
x
s
s x
+
\to -
-
=
=. 
\textcolor{lyred}{注}.
(
)
(
)
(
)
(
)
(
)
(
)
n
n
n
n
n
n
n
x
x
s
x
n n
n n
+
+
\infty 
\infty 
+
+
=
=
-
=
-
=
+
+
\sum 
\sum 
. 
\textcolor{lyred}{例}.求
-
+
-
+
-
+\Lambda 的和. 
解.考虑
(
)
n
n
n
x
n
\infty 
+
=
-
+
\sum 
,收敛域为[
]
1,1
-
,设()
(
)
n
n
n
s x
x
n
\infty 
+
=
-
=
+
\sum 
,则在(
)
1,1
-
上, 
()
(
)
(
)
n
n
n
n
n
s
x
x
x
x
\infty 
\infty 
=
=
'
=
-
=
-
=+
\sum 
\sum 
,于是()
()
arctan
x
s x
s
dx
x
x
-
=
=
+
\int 
,故 
(
)
()
n
n
s
n
\pi 
\infty 
=
-
=
=
+
\sum 
. 
\textcolor{lyred}{例}.求
(
)
n
n
n
\infty 
=
+
\sum 
的和. 
解.考虑
(
)
n
n
x
n
n
+
\infty 
=
+
\sum 
,收敛域为[
]
1,1
-
,设()
(
)
n
n
x
s x
n
n
+
\infty 
=
=
+
\sum 
,则在(
)
1,1
-
上, 
()
n
n
x
s
x
n
\infty 
=
'
=\sum 
,
()
n
n
n
n
x
x
s
x
x
x
x
x
x
x
\infty 
\infty 
-
=
=
''
=
=
=
\cdots 
=
-
-
\sum 
\sum 
,于是 
()
()
(
)
ln 1
x
x
s
x
s
dx
x
x
'
'
-
=
=-
-
-
\int 
, ()
()
(
)
ln 1
x
s x
s
x
dx
=
-
-
=
\int 
(
)
(
)
(
)
ln 1
2ln 1
x
x
x
x
-
-
+
-
+
,故
(
)
()
2ln 2
n
s
n
n
\infty 
=
=
=
-
+
\sum 
. 
\textcolor{lyred}{补充练习 }
1.填空题 
(1)设
n
n
n
a x
\infty 
=\sum 
有收敛半径3,则
(
)
n
n
n
na
x
\infty 
=
-
\sum 
的收敛区间为______. 
(2)设
(
)
n
n
n
a
x
\infty 
=
-
\sum 
在
x =
处发散,
x =-处收敛,则收敛域为______. 
(3)设
(
)
n
n
n
a
x
\infty 
=
+
\sum 
在
x =
处条件收敛,则收敛区间为______. 
解.(1)
R =
,(
)
2,4
-
;(2)
R =
,[
)
1,3
-
;(3)
R =
,(
)
5,3
-
. 
2.求
(
)
n
n
n
x
\infty 
=
+
\sum 
的和函数. 
解.令
t
x
=
,级数为
(
)
n
n
n
t
\infty 
=
+
\sum 
,收敛域为[
]
1,1
-
,设()
(
)
n
n
s t
n
t
\infty 
=
=
+
\sum 
,则 
在(
)
1,1
-
内, ()
(
)
(
)
n
n
n
n
t
s t
t
t
t
t
\infty 
\infty 
+
+
=
=
'
'




'
=
=
=
=




-


-


\sum 
\sum 
,故1
x
-<
<时, 
(
)
(
)
n
n
n
x
x
\infty 
=
+
=
-
\sum 
. 
3.求
n
n
n
x
n
\infty 
=
+
\sum 
的和函数. 
解.设()
n
n
n
s x
x
n
\infty 
=
+
=\sum 
,
(
)
1,1
x\in -
,则()
n
n
n
n
x
s x
nx
s
s
n
\infty 
\infty 
=
=
=
+
=
+
\sum 
\sum 
, 
()
(
)
n
n
n
n
x
x
s
x
x
nx
x
x
x
x
x
\infty 
\infty 
-
=
=
'
'




=
=
=
=




-


-


\sum 
\sum 
,
()
n
n
s
x
x
x
\infty 
-
=
'
=
=-
\sum 
,故 
()
()
(
)
ln 1
x
s
x
s
dx
x
x
-
=
=-
-
-
\int 
,于是()
(
)
(
)
ln 1
x
s x
x
x
=
-
-
-
. 
4.求
(
)(
)
n
n
x
n
n
\infty 
=
+
+
\sum 
的和函数. 
解.收敛域为[
]
1,1
-
,设()
(
)(
)
n
n
x
s x
n
n
\infty 
=
=
+
+
\sum 
,
(
)
1,1
x\in -
,则 
()
(
)(
)
()
n
n
x
x s x
s
x
n
n
+
\infty 
=
=
=
+
+
\sum 
,而
()
n
n
x
s
x
n
+
\infty 
=
'
=
+
\sum 
,
()
n
n
x
s
x
x
x
\infty 
=
''
=
=-
\sum 
, 
故
()
()
()
(
)
ln 1
x
x
s
x
s
x
s
dx
x
x
x
'
'
'
=
-
=
=-
-
-
-
\int 
,
()
()
1 0
s
x
s
-
= 
(
)
(
)
(
)
ln 1
ln 1
x
x
x
x dx
x
x
x
-
-
+
=
-
-
+
-




\int 
, ()
(
)
(
)
ln 1
x
x
x
s x
x
-
-
+
=
-
, 
()
()
lim
x
s
s x
\to 
=
=
, ()
()
lim
x
s
s x
-
\to 
=
=
. 
5.求
(
)
1 2
!
n
n
x
n
\infty 
=\sum 
的和函数. 
解.收敛域为(
)
,
-\infty +\infty ,设()
(
)
1 2
!
n
n
x
s x
n
\infty 
=
=\sum 
,则 
()
(
)
()
()
()
()
2 !
n
x
x
n
x
s
x
s x
s
x
s x
s x
C e
C e
n
-
\infty 
-
=
''
''
=
=
+\Rightarrow 
-
=\Rightarrow 
=
+
-
-
\sum 
, 
()
()
s
C
C
s
=

\Rightarrow 
=
=
'
=

,故
(
)
(
)
!
n
x
x
n
x
e
e
n
\infty 
-
=
=
+
-
\sum 
,
(
)
,
x\in -\infty +\infty . 
6.求
n
n
x
n
\infty 
=


+
+
+




\sum 
\Lambda 
的和函数. 
解.
lim 1
n
n
n
\to \infty 
+
+
+
=
\Lambda 
,收敛区间(
)
1,1
-
,设()
n
n
s x
x
n
\infty 
=


=
+
+
+




\sum 
\Lambda 
,则 
()
()
(
)
()
(
)
ln 1
ln 1
n
n
x
x
s x
xs x
x
s x
n
x
\infty 
=
-
-
=
=-
-
\Rightarrow 
=
-
\sum 
,
x <. 
7.求
2n
n
n
\infty 
=
+
\sum 
的和. 
解.考虑
(
)
n
n
n
x
\infty 
=
+
\sum 
,设()
(
)
n
n
s x
n
x
\infty 
=
=
+
\sum 
,
(
)
1,1
x\in -
,则()
s x = 
(
)
(
)
n
n
n
n
n
n
x
n
x
x
x
x
x
x
x
x
\infty 
\infty 
\infty 
+
=
=
=
'
'




+
-
=
-
=
-
=
-




-
-
-
-


-


\sum 
\sum 
\sum 
,于是 
n
n
n
s
\infty 
=
+


=
=
-
=




\sum 
. 
8.求
(
)
(
)
n
n
n
n
-
\infty 
=
-
-
\sum 
的和. 
解.考虑
(
)
(
)
n
n
n
x
n
n
-
\infty 
=
-
-
\sum 
,设()
(
)
(
)
n
n
n
s x
x
n
n
-
\infty 
=
-
=
-
\sum 
,
[
]
1,1
x\in -
,则 
()
(
)
n
n
n
s
x
x
n
-
\infty 
-
=
-
'
=
-
\sum 
,
()
(
)
n
n
s
x
x
x
\infty 
-
=
''
=
-
=+
\sum 
,故()
()
s x
s
'
'
-
= 
2arctan
x
dx
x
x
=
+
\int 
, ()
()
(
)
2 arctan
2 arctan
ln 1
x
s x
s
xdx
x
x
x
-
=
=
-
+
\int 
,于是 
()
ln 2
s
\pi 
=
-
. 
\vspace{6pt}
{\large \bfseries \textcolor{lyred}{第12.4 节 函数展开成幂级数}}

\vspace{4pt}
{\bfseries \textcolor{lyred}{一.Taylor 级数}}

在上一节,我们研究了如何求一个幂级数的和函数,现在讨论相反的问题,即如何 
把一个给定的函数
()
f x 展开为幂级数. 
若在I 内,
()
(
)
n
n
n
f x
a
x
x
\infty 
=
=
-
\sum 
,则称
()
f x 在I 内\textcolor{lyred}{可以展开为}
x
x
-
\textcolor{lyred}{的幂级数}. 
\textcolor{lyred}{定理(唯一性)}.设
()
f x 在
(
)
U x
内可以展开为
x
x
-
的幂级数
(
)
n
n
n
a
x
x
\infty 
=
-
\sum 
,则 
()(
)
!
n
n
f
x
a
n
=
,
0,1,2,
n =
\Lambda ,即
()
()(
)(
)
!
n
n
n
f
x
f x
x
x
n
\infty 
=
=
-
\sum 
. 
\textcolor{lyred}{定义}.称
()(
)(
)
!
n
n
n
f
x
x
x
n
\infty 
=
-
\sum 
为
()
f x 在
0x 处的\textcolor{lyred}{泰勒级数}; 
若它在
(
)
U x
内收敛,且以
()
f x 为和函数,称它为
()
f x 在
0x 处的\textcolor{lyred}{泰勒展开式}; 
特别地,
x =
时称
()()
!
n
n
n
f
x
n
\infty 
=\sum 
为
()
f x 的\textcolor{lyred}{麦克劳林级数}; 
若它在
()
U
内收敛,且以
()
f x 为和函数,则称它为
()
f x 的\textcolor{lyred}{麦克劳林展开式}. 
\textcolor{lyred}{定理}.设
()
f x 在
0x 处任意阶可导,则
()
f x 在
0x 邻的域内可展开为
x
x
-
的幂级数 
\Leftrightarrow 在该邻域内,
()
lim
n
n
R
x
\to \infty 
=
,其中
()
n
R
x 为
()
f x 的泰勒公式中的余项: 
()
(
)
(
)
(
)
(
)
(
)
1 !
n
n
n
f
x
x
x
R
x
x
x
n
\theta 
\theta 
+
+
+
-




=
-
<
<
+
. 
\vspace{4pt}
{\bfseries \textcolor{lyred}{二.函数展开成幂级数}}

{\bfseries \textcolor{lyred}{1.直接展开法}}

分两步:(1)写出
()()
!
n
n
n
f
x
n
\infty 
=\sum 
,求收敛半径R ;(2)在(
)
,
R R
-
内估计余项 
()
(
)(
)
(
)
(
)
1 0
1 !
n
n
n
f
x
R
x
x
n
\theta 
\theta 
+
+
=
<
<
+
,验证是否
()
lim
n
n
R
x
\to \infty 
=
. 
\textcolor{lyred}{例}.将
()
x
f x
e
=
展开为x 的幂级数. 
解.
()()
n
f
=,得
!
n
n
x
n
\infty 
=\sum 
, R =+\infty ,而
()
(
)
(
)
1 !
1 !
n
x
x
n
n
x
e
R
x
x
e
n
n
\theta 
+
+
=
\le 
+
+
,于是 
()
lim
n
n
R
x
\to \infty 
=
,故
!
n
x
n
x
e
n
\infty 
=
=\sum 
,
(
)
,
x\in -\infty +\infty . 
\textcolor{lyred}{例}.将
()
sin
f x
x
=
展开为x 的幂级数. 
解.
()()
sin 2
n
n
f
\pi 
=
,得
(
)(
)
1 !
n
n
n
x
n
+
\infty 
=
-
+
\sum 
, R =+\infty ,而 
()
(
)
(
)
sin
1 !
1 !
n
n
n
n
x
x
R
x
x
n
n
\pi 
\theta 
+
+


+




=
\le 
+
+
,于是
()
lim
n
n
R
x
\to \infty 
=
,故 
(
)(
)
sin
1 !
n
n
n
x
x
n
+
\infty 
=
=
-
+
\sum 
,
(
)
,
x\in -\infty +\infty . 
\textcolor{lyred}{例(牛顿二项展开式)}.将
()
(
)
f x
x
\alpha 
=
+
展开成x 的幂级数. 
解.不妨设n\notin \infty ,
()()
(
)
(
)(
)
1 1
n
n
f
x
n
x
\alpha 
\alpha \alpha 
\alpha 
-
=
-
-
+
+
\Lambda 
,得
()
f
=, 
()
f
\alpha 
'
=
,
()
(
)
f
\alpha \alpha 
''
=
-
,\Lambda ,
()()
(
)
(
)
n
f
n
\alpha \alpha 
\alpha 
=
-
-
+
\Lambda 
,得麦克劳林级数 
(
)
(
)
(
)
1!
2!
!
n
n
n
n
x
x
x
x
n
n
\alpha 
\alpha \alpha 
\alpha \alpha 
\alpha 
\alpha 
\infty 
=
-
-
-
+


+
\cdots +
+
+
+
=
\cdots 




\sum 
\Lambda 
\Lambda 
\Lambda 
,其中 
(
)
(
)
!
n
n
n
\alpha 
\alpha \alpha 
\alpha 
-
-
+

=




\Lambda 
,
lim
lim
n
n
n
n
a
n
R
a
n
\alpha 
\rho 
+
\to \infty 
\to \infty 
-
=
=
=\Rightarrow 
=
+
; 
设()
n
n
s x
x
n
\alpha 
\infty 
=


=
\cdots 




\sum 
,
(
)
1,1
x\in -
,则(
)()
()
()
(
)
x s x
s x
s x
C
x
\alpha 
\alpha 
'
+
=
\Rightarrow 
=
+
,且 
()
s
C
=\Rightarrow 
=,故(
)
n
n
x
x
n
\alpha 
\alpha 
\infty 
=


+
=
\cdots 




\sum 
,
(
)
1,1
x\in -
. 
{\bfseries \textcolor{lyred}{2.间接展开法}}

\textcolor{lyred}{例}.将
()
f x
x
x
=
+
+
展开为x 的幂级数. 
解.
()
(
)(
)
(
)
(
)
2 1
6 1
x
f x
x
x
x
x
x


=
=
-
=
\cdots 
-
\cdots 
=


+
+
+
+
--
--


(
)
(
)
(
)
n
n
n
n
n
n
n
n
n
n
x
x
x
\infty 
\infty 
\infty 
+
=
=
=




-
-
-
=
-
-








\sum 
\sum 
\sum 
, 1
x
-<
<. 
\textcolor{lyred}{例}.将
()
f x
x
x
=
+
+
展开为
x -的幂级数. 
解.设
t
x
=
-,则
()
(
)(
)
(
)(
)
f x
x
x
t
t
t
t


=
=
=
-
=


+
+
+
+
+
+


(
)
(
)
(
)
4 1
8 1
n
n
n
n
n
n
t
t
t
t
t
\infty 
\infty 
=
=




\cdots 
-
\cdots 
=
-
-
-
-<<
=




+
+




\sum 
\sum 
(
)
(
)
(
)
n
n
n
n
n
n
n
n
n
n
t
x
\infty 
\infty 
+
+
+
+
=
=




-
-
=
-
-
-








\sum 
\sum 
, 1
x
-<
<. 
\textcolor{lyred}{注}.
n
n
u
u
\infty 
=
=
-
\sum 
,
(
)
n
n
n
u
u
\infty 
=
=
-
+
\sum 
,其中
u <. 
\textcolor{lyred}{例}.将
()
(
)
f x
x
=
-
展开为x 的幂级数. 
解.
()
2 1
n
n
n
n
n
x
x
n
f x
x
x
\infty 
\infty 
-
+
=
=
'
'
'








=
=
=
=










-
-










\sum 
\sum 
, 2
x
-<
<
. 
\textcolor{lyred}{例}.将
()
cos
f x
x
=
展开为x 的幂级数. 
解.
()
(
)
(
)(
)
(
)(
)
sin
1 !
!
n
n
n
n
n
n
x
x
f x
x
n
n
+
\infty 
\infty 
=
=
'


'
=
=
-
=
-


+


\sum 
\sum 
,
(
)
,
x\in -\infty +\infty . 
\textcolor{lyred}{例}.将
()
(
)
ln 1
f x
x
=
+
展开为x 的幂级数. 
解.
()
(
)
(
)
(
)
n
n
n
n
n
f
x
x
x
x
x
\infty 
\infty 
=
=
'
=
=
=
-
=
-
+
--
\sum 
\sum 
, 1
x
-<
<,故 
()
()
()
(
)
x
n
n
n
x
f x
f
f
x dx
n
+
\infty 
=
'
=
+
=
-
+
\sum 
\int 
,
(
]
1,1
x\in -
. 
\textcolor{lyred}{例}.将
()
(
)
(
)
ln 1
f x
x
x
=
-
+
展开为x 的幂级数. 
解一.
()
(
)
(
)
(
)
ln 1
n
n
n
n
n
n
f
x
x
x
x
x
n
\infty 
\infty 
+
=
=
'
=-+
-
+
=-+
-
-
-
=
+
+
\sum 
\sum 
(
)
(
)
(
)
n
n
n
n
n
n
n
n
n
n
x
x
x
n
n
\infty 
\infty 
\infty 
-
=
=
=
+
+
-
-
-
=+
-
\sum 
\sum 
\sum 
, 1
x
-<
<,故 
()
()
()
(
)
(
)
x
n
n
n
n
f x
f
f
x dx
x
x
n n
\infty 
+
=
+
'
=
+
=
+
-
+
\sum 
\int 
,
(
]
1,1
x\in -
. 
解二.
()
(
)
(
)
(
)
(
)
n
n
n
n
n
n
n
n
n
x
f x
x
x
x
n
n
n
+
\infty 
\infty 
\infty 
+
+
=
=
=
=
-
-
=
-
-
-
=
+
+
+
\sum 
\sum 
\sum 
(
)
(
)
(
)
(
)
n
n
n
n
n
n
n
n
n
n
x
x
x
x
x
n
n
n n
\infty 
\infty 
\infty 
-
+
+
+
=
=
=
+
+
-
-
-
=
+
-
+
+
\sum 
\sum 
\sum 
,
(
]
1,1
x\in -
. 
\textcolor{lyred}{例}.将
()
arctan
f x
x
=
展开为x 的幂级数. 
解.
()
(
)
(
)
(
)
n
n
n
n
n
f
x
x
x
x
x
\infty 
\infty 
=
=
'
=
=
=
-
=
-
+
--
\sum 
\sum 
, 1
x
-<
<,故 
()
()
()
(
)
x
n
n
n
x
f x
f
f
x dx
n
+
\infty 
=
'
=
+
=
-
+
\sum 
\int 
,
[
]
1,1
x\in -
. 
\textcolor{lyred}{例}.将
()
arctan 1
x
f x
x
+
=
-
展开为x 的幂级数. 
解.
()
(
)
(
)
(
)
n
n
n
n
n
f
x
x
x
x
x
\infty 
\infty 
=
=
'
=
=
=
-
=
-
+
--
\sum 
\sum 
,故 
()
()
()
(
)
x
n
n
n
x
f x
f
f
x dx
n
\pi 
+
\infty 
=
'
=
+
=
+
-
+
\sum 
\int 
,
[
)
1,1
x\in -
. 
\textcolor{lyred}{例}.设
()
sin
,
1,
x
x
f x
x
x

\neq 

=

=

,求
()()
n
f
,
1,2,3,
n =
\Lambda . 
解.当
x \neq 
时,
()
(
)(
)
sin
1 !
n
n
n
x
x
f x
x
n
\infty 
=
=
=
-
+
\sum 
,由连续性,当
x =
时也成立,故 
由展开的唯一性,
(
)()
1 0
n
f
-
=
,
(
)()
(
)
(
)(
)
(
)
!
1 !
n
n
n
f
n
n
n
-
-
=
\cdots 
=
+
+. 
\textcolor{lyred}{补充练习 }
1.求
!
n
n
n
x
n
\infty 
=
+
\sum 
的和函数. 
解. ()
(
)
!
!
1 !
!
n
n
n
n
x
x
n
n
n
n
n
x
x
s x
x
x
x
xe
e
n
n
n
n
-
\infty 
\infty 
\infty 
\infty 
=
=
=
=
=
+
=
+
=
+
-
-
\sum 
\sum 
\sum 
\sum 
,
x
-\infty <
<\infty ; 
或者, ()
(
)
!
!
!
n
n
n
x
x
x
n
n
n
n
x
x
s x
x
x
x e
xe
e
n
n
n
+
\infty 
\infty 
\infty 
=
=
=
'
'




+
'


=
=
=
=
-
=
+
-










\sum 
\sum 
\sum 
. 
2.将
()
x
f
x
x
x
=+
-
展开为x 的幂级数. 
解.
()
(
)(
)
(
)(
)
3 1
x
f x
x
x
x
x
x
x
x
+
-


=
=
-
=
-
=


+
-
+
-
+
-
+


(
)
(
)
n
n
n
n
n
n
n
x
x
x
\infty 
\infty 
\infty 
=
=
=
--


-
-
=




\sum 
\sum 
\sum 
,
x
-
<
<
. 
3.将
()
ln1
x
f x
x
=
+
展开为
x -的幂级数. 
解.
()
(
)
(
)
n
n
n
n
x
f
x
x
x
x
x
x
\infty 
\infty 
=
=
-


'
=
-
=
-
\cdots 
=
-
-
-
=


-
+
-
-


+
\sum 
\sum 
(
)
(
)
n
n
n
n
x
\infty 
+
=


-
-
-




\sum 
,故
()
()
()
x
f x
f
f
x dx
'
=
+
=
\int 
(
)
(
)
ln
n
n
n
n
x
n
\infty 
+
+
=
-


+
-
-


+


\sum 
,
x
x
-
<\Leftrightarrow 
<
<
. 
4.设
()
5 sin
f x
x
x
=
,求
(
)()
f
. 
解.
()
3!
5!
3!
5!
x
x
x
x
f x
x
x
x


=
-
+
-
=
-
+
-




\Lambda 
\Lambda ,故
(
)()
10!
5!
f
=
. 
\vspace{6pt}
{\large \bfseries \textcolor{lyred}{第12.7 节 傅里叶级数}}

\vspace{4pt}
{\bfseries \textcolor{lyred}{一.三角函数系的正交性}}

\textcolor{lyred}{定义}.称\{
\}
1,cos ,sin ,cos2 ,sin 2 ,
,cos
,sin
,
x
x
x
x
nx
nx
\Lambda 
\Lambda 为\textcolor{lyred}{三角函数系}.\textcolor{lyred}{ }
\textcolor{lyred}{定理}.三角函数系在[
]
,
\pi \pi 
-
上\textcolor{lyred}{正交},即其中任何两个不同函数的积在[
]
,
\pi \pi 
-
上的 
积分等于零. 
\textcolor{lyred}{定义}.称
cos
sin
n
n
n
a
a
nx
b
nx
\infty 
=
+
+
\sum 
为以2\pi 为周期的\textcolor{lyred}{三角级数};一般地, 
称
cos
sin
n
n
n
a
n x
n x
a
b
l
l
\pi 
\pi 
\infty 
=


+
+




\sum 
为以2l 为周期的\textcolor{lyred}{三角级数}. 
\vspace{4pt}
{\bfseries \textcolor{lyred}{二.函数展开成傅里叶级数}}

\textcolor{lyred}{定义}.记
()
(
)
cos
na
f x
nxdx n
\pi 
\pi 
\pi -
=
\ge 
\int 
,
()
(
)
sin
nb
f x
nxdx n
\pi 
\pi 
\pi -
=
\ge 
\int 
,称为
()
f x 的 
\textcolor{lyred}{傅里叶系数},此时三角级数
(
)
cos
sin
n
n
n
a
a
nx
b
nx
\infty 
=
+
+
\sum 
称为
()
f x 的\textcolor{lyred}{傅里叶级数}, 
部分和
()
n
F
x 称为
()
f x 的\textcolor{lyred}{傅里叶多项式}. 
\vspace{4pt}
{\bfseries \textcolor{lyred}{三.收敛的一个充分条件}}

\textcolor{lyred}{定理(收敛定理---狄利克雷充分条件)}.设
()
f x 具有周期2\pi ,若它在一个周期内 
满足(1)连续,或者只有有限多个第一类间断点;(2)只有至多有限多个极值点,则 
()
f x 的傅里叶级数收敛,其和函数为 
()
(
)
(
)
(
)
cos
sin
n
n
n
a
f x
a
nx
b
nx
f x
f x
\infty 
-
+
=


=
+
+
=
+


\sum 

; 
特别地,在连续点处,
()
(
)
cos
sin
n
n
n
a
f x
a
nx
b
nx
\infty 
=
=
+
+
\sum 
. 
\textcolor{lyred}{例}.设
()
f x 具有周期2\pi ,且
()
1,
1,
x
f x
x
\pi 
\pi 
-
-
\le 
<

=
\le 
<

,将
()
f x 展开为傅里叶级数. 
解.
()
(
)
cos
1 cos
1 cos
na
f x
nxdx
nxdx
nxdx
\pi 
\pi 
\pi 
\pi 
\pi 
\pi 
\pi 
-
-
=
=
-
+
\cdots 
=
\int 
\int 
\int 
, 
()
(
)
sin
1 sin
1 sin
sin
nb
f x
nxdx
nxdx
nxdx
nxdx
\pi 
\pi 
\pi 
\pi 
\pi 
\pi 
\pi 
\pi 
\pi 
-
-
=
=
-
+
\cdots 
=
=
\int 
\int 
\int 
\int 
(
)
cos
2 1 cos
n
nx
n
n
n
n
\pi 
\pi 
\pi 
\pi 
\pi 


--
-




-
=
\cdots 
=
\cdots 




,故 
()
(
)
sin
sin3
sin 2
f x
x
x
k
x
k
\pi 


=
+
+
+
-
+


-


\Lambda 
\Lambda ,其中x
k\pi 
\neq 
. 
\textcolor{lyred}{注}.在x
k\pi 
=
处,
()
1 1
f x
-+
=
=

. 
\textcolor{lyred}{例}.设
()
f x 具有周期2\pi ,且
()
,
0,
x
x
f x
x
\pi 
\pi 
-
\le 
<

=
\le 
<

,将
()
f x 展开为傅里叶级数. 
解.
()
cos
cos
sin
na
f x
nxdx
x
nxdx
xd
nx
n
\pi 
\pi 
\pi 
\pi 
\pi 
\pi 
\pi 
-
-
-
=
=
=
=
\int 
\int 
\int 
(
)
cos
1 cos
sin
n
nx
n
nxdx
n
n
n
n
n
\pi 
\pi 
\pi 
\pi 
\pi 
\pi 
\pi 
-
-


--
-


-
=
=
=








\int 
, 
()
a
f x dx
xdx
\pi 
\pi 
\pi 
\pi 
\pi 
\pi 
-
-
=
=
=-
\int 
\int 
,
()
sin
sin
nb
f x
nxdx
x
nxdx
\pi 
\pi 
\pi 
\pi 
\pi 
-
-
=
=
=
\int 
\int 
(
)
cos
cos
cos
cos
n
n
xd
nx
n
nxdx
n
n
n
n
\pi 
\pi 
\pi 
\pi 
\pi 
\pi 
\pi 
-
-
-


-
-
=-
-
=-
=




\int 
\int 
,故 
()
cos
sin
sin 2
cos3
sin3
sin 4
f x
x
x
x
x
x
x
\pi 
\pi 
\pi 




=-
+
+
-
+
+
-
+








cos5
sin5
sin 6
x
x
x
\pi 


+
-
+




\Lambda ,其中
(
)
x
k
\pi 
\neq 
+
. 
\textcolor{lyred}{例}.设
()
,
,
x
x
f x
x
x
\pi 
\pi 
-
-
\le 
<

=
\le 
\le 

,将
()
f x 展开为具有周期2\pi 的傅里叶级数. 
解.将
()
f x 延拓为(
)
,
-\infty +\infty 上周期为2\pi 的函数,它是连续偶函数,则 
()
a
f x dx
xdx
\pi 
\pi 
\pi 
\pi 
\pi 
\pi 
-
=
=
=
\int 
\int 
,
()
sin
nb
f x
nxdx
\pi 
\pi 
\pi -
=
=
\int 
, 
()
()
cos
cos
cos
na
f x
nxdx
f x
nxdx
x
nxdx
\pi 
\pi 
\pi 
\pi 
\pi 
\pi 
\pi 
-
=
=
=
=
\int 
\int 
\int 
(
)
cos
sin
sin
n
n
xd
nx
nxdx
n
n
n
n
n
\pi 
\pi 
\pi 
\pi 
\pi 
\pi 
\pi 


-
-
-
=
-
=
\cdots 
=
\cdots 




\int 
\int 
,故 
()
cos
cos3
cos5
f x
x
x
x
\pi 
\pi 


=
-
+
+
+




\Lambda ,
x
\pi 
\pi 
-
\le 
\le 
. 
\textcolor{lyred}{注}.令
x =
,则
\pi 
\pi 
\pi 


=
-
+
+
+
\Rightarrow 
=+
+
+




\Lambda 
\Lambda . 
\textcolor{lyred}{例}.求和
\sigma =+
+
+
+\Lambda 与
\tau =-
+
-
+\Lambda . 
解.记
\pi 
\sigma =
+
+
+
=
\Lambda 
,
\sigma 
\sigma 
=
+
+
+
=
\Lambda 
,则 
\pi 
\pi 
\sigma 
\sigma 
\sigma 
\sigma 
\sigma 
=
+
=
+
\Rightarrow 
=
,而
\pi 
\sigma 
\sigma 
=
=
,故 
\pi 
\tau 
\sigma 
\sigma 
=
-
=
. 
\vspace{4pt}
{\bfseries \textcolor{lyred}{四.正弦级数和余弦级数}}

若()
f
x 是具有周期2\pi 的奇函数,则
na =
,
()
sin
nb
f
x
nxdx
\pi 
\pi 
=
\int 
,此时,它的 
傅立叶级数为
sin
n
n
b
nx
\infty 
=\sum 
,称为\textcolor{lyred}{正弦级数}; 
若()
f
x 是具有周期2\pi 的偶函数,则
()
cos
na
f
x
nxdx
\pi 
\pi 
=
\int 
,
nb =
,此时,它的 
傅立叶级数为
cos
n
n
a
a
nx
\infty 
=
+\sum 
,称为\textcolor{lyred}{余弦级数}. 
\textcolor{lyred}{例}.设
()
f x 具有周期2\pi ,且在[
)
,
\pi \pi 
-
上
()
f x
x
=
,将
()
f x 展开为傅里叶级数. 
解.若不计
(
)
x
k
\pi 
=
+
(对定积分没有影响),则
()
f x 是奇函数,故 
na =
,
()
sin
sin
cos
nb
f x
nxdx
x
nxdx
xd
nx
n
\pi 
\pi 
\pi 
\pi 
\pi 
\pi 
=
=
=-
=
\int 
\int 
\int 
(
)
cos
cos
cos
n
n
nxdx
n
n
n
n
\pi 
\pi 
\pi 
\pi 
\pi 
-


-
+
=-
=
-




\int 
,故 
()
2 sin
sin 2
sin3
f x
x
x
x


=
-
+
-




\Lambda ,
(
)
x
k
\pi 
\neq 
+
. 
\textcolor{lyred}{注}.在
(
)
x
k
\pi 
=
+
处,
()
f x
\pi 
\pi 
-
=
=

. 
\textcolor{lyred}{例}.将
()
(
)
sin 2
x
f x
E
x
\pi 
\pi 
=
-
\le 
\le 
展开为具有周期2\pi 的傅里叶级数,其中
E >
. 
解.将
()
f x 延拓为(
)
,
-\infty +\infty 上周期为2\pi 的函数,它是连续偶函数,则 
()
(
)
cos
sin
cos
n
x
E
a
f x
nxdx
E
nxdt
n
\pi 
\pi 
\pi 
\pi 
\pi 
=
=
=-
-
\int 
\int 
,
nb =
,故 
()
cos
n
E
E
f x
nx
n
\pi 
\pi 
\infty 
=
=
-
-
\sum 
,
x
\pi 
\pi 
-
\le 
\le 
. 
\textcolor{lyred}{例}.将
()
(
)
x
f x
x
\pi 
\pi 
-
=
\le 
\le 
分别展开为具有周期2\pi 的正弦级数和余弦级数. 
解.奇延拓到[
]
,
\pi \pi 
-
上,
()
sin
sin
nb
f x
nxdx
nxdx
n
\pi 
\pi \pi 
\pi 
\pi 
-
=
=
=
\int 
\int 
,故 
()
1 sin
n
f x
nx
n
\infty 
=
=\sum 
,
(
]
0,
x
\pi 
\in 
; 
偶延拓到[
]
,
\pi \pi 
-
上,
()
x
a
f x dx
dx
\pi 
\pi \pi 
\pi 
\pi 
\pi 
-
=
=
=
\int 
\int 
, 
()
(
)
1 cos
cos
cos
n
n
x
n
a
f x
nxdx
nxdx
n
n
\pi 
\pi \pi 
\pi 
\pi 
\pi 
\pi 
\pi 
--
-
-
=
=
=
=
\int 
\int 
,故 
()
(
)
(
)
cos 2
n
f x
n
n
\pi 
\pi 
\pi 
\infty 
=
=
+
-
-
\sum 
,0
x
\pi 
\le 
\le 
. 
\textcolor{lyred}{补充练习 }
1.设
()
f x 具有周期2\pi ,且
()
,
1,
x
x
f x
x
\pi 
\pi 
-
\le 
<

=
\le 
<

,它的傅里叶系数为
0a ,
na ,
nb , 
求
n
n
a
a
\infty 
=
+\sum 
. 
解.
()
(
)
(
)
0 1
n
n
a
a
f
f
f
\infty 
-
+
=
+


+
=
=
+
=
=


\sum 

. 
2.将()
(
)
1 0
f
x
x
x
\pi 
=
+
\le 
\le 
分别展开为具有周期2\pi 的正弦级数和余弦级数. 
解.奇延拓到[
]
,
\pi \pi 
-
上,
()
(
)
sin
1 sin
nb
f x
nxdx
x
nxdx
\pi 
\pi 
\pi 
\pi 
=
=
+
=
\int 
\int 
(
)
(
)(
)
1 cos
n
n
n
n
\pi 
\pi 
\pi 
\pi 
\pi 


-
+
-
=-
+
-
-





,故 
(
)
(
)
2 sin
sin 2
2 sin 3
sin 4
x
x
x
x
x
\pi 
\pi 
\pi 
\pi 
\pi 
+
-
+
+
-
+


+=




\Lambda ,0
x
\pi 
<
<
; 
偶延拓到[
]
,
\pi \pi 
-
上,
()
(
)
cos
1 cos
na
f x
nxdx
x
nxdx
\pi 
\pi 
\pi 
\pi 
=
=
+
=
\int 
\int 
(
)
(
)
cos
sin
sin
n
nx
x
d
nx
nxdx
n
n
n
n
n
n
\pi 
\pi 
\pi 
\pi 
\pi 
\pi 
-
-
-
+
=-
=
\cdots 
=
\cdots 
\int 
\int 
, 
(
)
a
x
dx
\pi 
\pi 
\pi 
=
+
=
+
\int 
,故 
()
cos
cos3
cos5
f x
x
x
x
\pi 
\pi 


=
+-
+
+
+




\Lambda ,0
x
\pi 
\le 
\le 
. 
\vspace{6pt}
{\large \bfseries \textcolor{lyred}{第12.8 节 一般周期函数的傅里叶级数}}

\textcolor{lyred}{定理}.设
()
f x 是
T
l
=
的周期函数,且满足收敛定理的条件,则在连续点处 
()
cos
sin
n
n
n
a
n x
n x
f x
a
b
l
l
\pi 
\pi 
\infty 
=


=
+
+




\sum 
,其中 
()
(
)
cos
0,1,2,
l
n
l
n x
a
f x
dx n
l
l
\pi 
-
=
=
\int 
\Lambda ,
()
(
)
sin
1,2,
l
n
l
n x
b
f x
dx n
l
l
\pi 
-
=
=
\int 
\Lambda ; 
而在间断点处,该级数收敛到
(
)
(
)
f x
f x
-
+
+
. 
\textcolor{lyred}{注}.上述积分区间可以换为任何长2l 的区间,例如[
]
0,2l . 
\textcolor{lyred}{注}.当
()
f x 为奇函数时,
na =
,
()
sin
l
n
n x
b
f x
dx
l
l
\pi 
=\int 
; 
当
()
f x 为偶函数时,
()
cos
l
n
n x
a
f x
dx
l
l
\pi 
=\int 
,
nb =
. 
\textcolor{lyred}{例}.设
()
f x 具有周期4 ,并且
()
0,
1,
x
f x
x
-\le 
<

=
\le 
<

,将它展开为傅里叶级数. 
解.
()
a
f x dx
dx
-
=
=
=
\int 
\int 
,
()
cos
cos
n
n x
n x
a
f x
dx
dx
\pi 
\pi 
-
=
=
=
\int 
\int 
, 
()
(
)
(
)
sin
sin
1 cos
n
n
n x
n x
b
f x
dx
dx
n
n
n
\pi 
\pi 
\pi 
\pi 
\pi 
-


=
=
=
-
=
--


\int 
\int 
, 
故
()
sin
sin
sin
x
x
x
f x
\pi 
\pi 
\pi 
\pi 


=
+
+
+
+




\Lambda ,
x
k
\neq 
. 
\textcolor{lyred}{例}.将
()
(
)
2 0
f x
x
x
=
\le 
\le 
分别展开为正弦级数和余弦级数. 
解.奇延拓到[
]
2,2
-
上,得到
()
sin
n
n
n x
f
x
b
\pi 
\infty 
=
=\sum 
,0
x
\le 
<
,其中 
()
(
)
(
)
sin
sin
n
n
n
n x
n x
b
f x
dx
x
dx
n
n
\pi 
\pi 
\pi 
\pi 
+


=
=
=-
+
-
-


\int 
\int 
; 
偶延拓到[
]
2,2
-
上,得到
()
cos
n
n
a
n x
f x
a
\pi 
\infty 
=
=
+\sum 
,0
x
\le 
\le 
,其中 
()
a
f x dx
x dx
=
=
=
\int 
\int 
,
()
(
)
cos
cos
n
n
n x
n x
a
f x
dx
x
dx
n
\pi 
\pi 
\pi 
=
=
=-
\int 
\int 
. 
\textcolor{lyred}{补充练习 }
1.设
()
f x 具有周期2 ,且
()
1,
,
x
f x
x
x

-\le 
\le 

=

<
<

,则其傅里叶级数在
x =-
处的 
和为______.\textcolor{lyred}{ }
解.
s
s
s








-
=
-
+
=
=
+
=
















. 
2.设
()
f x
x
=
+,0
x
\le 
<,而()
sin
n
n
s x
b
n x
\pi 
\infty 
=
=\sum 
,
(
)
,
x\in -\infty \infty ,其中 
()
(
)
sin
1,2,
nb
f x
n xdx n
\pi 
=
=
\int 
\Lambda ,则
s

-
=




______. 
解.
s
s
f






-
=-
=-
=-












. 
3.设
()
,
2 ,
x
x
f x
x
x

\le 
\le 

=
-
<
<

,而()
cos
n
n
a
s x
a
n x
\pi 
\infty 
=
=
+\sum 
,
(
)
,
x\in -\infty \infty ,其中 
()
cos
na
f x
n xdx
\pi 
=\int 
,则
s

-
=




______. 
解.
1 1
2 2
s
s
s
s










-
=
-
+
=
-
=
=
+
=




















. 